% Pursuing Stacks (cotes 134-2 to 134-8): archive cote 134-7, site page 70 (PDF page 71), typescript page 512.
% typescript. Draft compared: ps6.tex:3468-3545.
% AI-transcribed by Claude Opus 5.5 (claude-opus-5-5, adaptive thinking, effort high), Message Batches API, batch ps4 (msgbatch_012Aoc1NuskATPsyBeKSYK7J), 2026-11-12; not yet checked by eye.
\scanpage{134-7/70}
\mspage{512}

Moreover the canonical map
\[ X \longrightarrow W(X) = \underline{Z}^{(X)} \]
factors (by construction) through
\[ (3) \qquad X \longrightarrow W(X)(1) \quad ( \hookrightarrow W(X) ) , \]
and this map is universal for all maps of X into torsors under abelian
group objects of $\Delta^{\wedge}$ . It is an immediate consequence of Whiteheads
theorem tha\over{r}{t} a weak equivalence $X \to Y$ in $\Delta^{\wedge}$ induces \xout{w} weak equivalen-%
ce\hand{s} in $\Delta^{\wedge}$
\[ (4) \qquad \hand{\widetilde{W}}(X) \longrightarrow \hand{\widetilde{W}}(Y) \quad , \qquad W(X)(1) \longrightarrow W(Y)(1) \]
(and also $W(X)(n) \to W(Y)(n)$ for any $n \in \underline{Z}$ ), and accordingly that
the map between the normalized chain complexes corresponding to $\hand{\widetilde{W}}(X)$ and
$\hand{\widetilde{W}}(Y)$, namely by definition
\[ (4') \qquad L\hand{\widetilde{H}}.(X) \longrightarrow L\hand{\widetilde{H}}.(Y) \]
is a quasi-isomorphism. Thus we get \xout{functors} a functor
\[ \hand{\widetilde{LH}}. : \mathrm{Hot} \longrightarrow \mathrm{Hotab} \quad , \]
whose composition\inserted{$L\widetilde{H}.\circ\beta$ ,} with the forgetful functor
\[ \hand{\beta} : \mathrm{Hot}^{\bullet} \longrightarrow \mathrm{Hot} \]
is \xout{thexfunct} canonically isomorphic to the functor $L\hand{\widetilde{H}}.$ of p.510 (formula
9). More specifically, when X is in $\Delta^{\wedge\bullet}$ , then the map
\[ \underline{Z} = \underline{Z}^{(e)} \longrightarrow \underline{Z}^{(X)} \quad \text{deduced from } e \to X \]
defines a splitting of the extension 1, i.e. an isomorphism
\[ (5) \qquad \underline{Z}^{(X)} = W(X) \simeq \hand{\widetilde{W}}(X) \oplus \underline{Z} \quad , \]
and accordingly $\hand{\widetilde{W}}(X)$ may be interpreted, at will, as a quotient group
of W(X) (which we did on p.510)\hand{,} or as a subobject of W(X),\deleted{(}which is the
better choice, because it makes a good sense without using any pointed
structure. Again, without using the pointed structure, we\xout{x} get a canoni-%
cal exact trian\xout{x}gle in D(Ab) interpreting (1)
\diagram{
\begin{tikzcd}[column sep=small]
 & \underline{Z} \arrow[dl] & \\
L\hand{\widetilde{H}}.(X) \arrow[rr] & & LH.(X) \arrow[ul]
\end{tikzcd}
}
(6)

replacing ((12) p. 510), which \deleted{should be replaced} (in the case considered
there, namely X pointed\hand{)} \inserted{should be replaced} by the more precise relationship
\[ (\xout{4}\, 5') \qquad LH.(X) \simeq L\hand{\widetilde{H}}.(X) \oplus \underline{Z} \quad , \]
an isomorphism functorial for X in Hot\hand{$^{\bullet}$}.

\margin{70}
